\chapter{Schwingungen und Wellen}

Schwingungen und Wellen sind zentrale Phänomene in der Physik, die eng miteinander 
verknüpft sind und in vielen technischen Anwendungen eine entscheidende Rolle spielen. 

Schwingungen beschreiben die periodische Bewegung eines Systems um eine 
Gleichgewichtslage - wie das Pendel einer Uhr oder die Vibrationen einer Maschine. 
Sie sind die Grundlage für das Verständnis von Wellen, die sich als räumliche 
Ausbreitung von Schwingungen beschreiben lassen, ohne dass Material selbst wandert.

\section{Schwingungen}

\begin{minipage}{.68\textwidth}
    Die Untersuchung von Schwingungen ist eine wichtige Aufgabe der Physik, da
    Schwingungen an vielen technischen Vorgängen beteiligt sind. Sie treten als 
    Vibrationen von Bauwerken und Maschinenfundamenten ebenso auf wie beim Einschwingen 
    von Messgeräten und Regeleinrichtungen in der Automatisierungstechik. Schwingende Pendel, 
    Unruhen und Schwingquarze sind Grundlage der meisten Zeitmessverfahren.

    In diesem Abschnitt werden die Grundlagen von Schwingungen behandelt, 
    um in einem späteren Schritt Wellen besser verstehen zu können. 
\end{minipage}
\hspace{.05\textwidth}
\begin{minipage}{.35\textwidth}
    \begin{figure}[H]   
        \centering          %     left botm right top
        \includegraphics[clip, trim=0cm 0cm 0cm 0cm,width=\textwidth]{pendel_bild.jpg} %.6 stands for 0.6 x textwidth
        \caption{Klassisches Beispiel für eine Schwingung: Fadenpendel \cite{noauthor_vergleich_nodate}}
        \label{fig:pendel}
    \end{figure}
\end{minipage}

Abbildung \ref{tab:begriffeundbezeichnungen} zeigt Begriffe und Grössen, wobei Unterschiede zu den
in der Technik üblichen Bezeichnungen nicht zu vermeiden sind.

\begin{table}[H]
    \centering
    \caption{Begriffe und Bezeichnungen \cite{bartsch_dynamik_2023}}
    \label{tab:begriffeundbezeichnungen}
    \begin{figure}[H]    
        \centering          %     left botm right top
        \includegraphics[width=.7\textwidth]{begriffe_schwingungen.png} %.6 stands for 0.6 x textwidth
    \end{figure}
\end{table}

\subsection{Freie, ungedämpfte Schwingung}

Das Fadenpendel wie in Abbildung \ref{fig:pendel} dargestellt ist ein Beispiel 
für die freie, ungedämpfte Schwingung (wenn der Widerstand vernachlässigt wird).

Das Federpendel verhaltet sich ähnlich, die auftretende Bewegung folgt einer Sinusschwingung. 

Wir wollen am typischen Beispiel des Masse-Feder-Schwingers die Bewegung eines
Körpers beschreiben, der an einer Schraubenfeder hängt und freie ungedämpfte
Schwingungen um seine Gleichgewichtslage ausführt. Der Körper lässt sich dabei
als Punktmasse auffassen, die eine spezielle geradlinige, ungleichmässig beschleunigte 
Bewegung ausführt \cite{zeitler_physik_2016}.

\subsubsection{Kinematik der Sinusschwingung}

In allgemeiner Form können Auslenkung $x$, 
Geschwindigkeit $v$ und Beschleunigung $a$ wie folgt ausgedrückt werden:

\begin{equbox}
    \begin{align}
        x(t) &= A \cdot \sin(\omega_0 \cdot t + \phi) \\
        v(t) &= A \cdot \omega_0 \cdot \cos(\omega_0 \cdot t + \phi) \\
        a(t) &= - A \cdot \omega_0^2 \cdot \sin(\omega_0 \cdot t + \phi)
    \end{align}
\end{equbox}

Die Amplitude $A$ und die Phasenverschiebung $\phi$ sind aus den Anfangsbedingungen, d.h.
Position und Geschwindigkeit bei $t = 0$ zu bestimmen:

\begin{equbox}
    \begin{align}
        A &= \sqrt{x_0^2 + (\frac{v_0}{\omega_0})^2} \\
        \phi &= \tan^{-1}(\frac{x_0 \cdot \omega_0}{v_0})
    \end{align}
\end{equbox}

Beim Phasenwinkel ist auf die Mehrdeutigkeit der arctan-Funktion zu achten!

Folgend die Umsetzung an einem Massenschwinger. 
Die Auslenkung $y$ aus der Nulllage ist in Abhängigkeit von der Zeit
in Bild \ref{fig:freie_ungedämpfte} a) dargestellt. Das y-t-Diagramm zeigt eine Sinuskurve.

\begin{figure}[H]   
    \centering          %     left botm right top
    \includegraphics[clip, trim=0cm 0cm 0cm 0cm,width=.5\textwidth]{freie_ungedämpfte_schwingung.png} %.6 stands for 0.6 x textwidth
    \caption{Freie ungedämpfte Sinusschwingung a) y-t-Diagramm; b) v-t-Diagramm; c) a-t-Diagramm \cite{zeitler_physik_2016}}
    \label{fig:freie_ungedämpfte}
\end{figure}

Mit der Amplitude $y_{max}$ als maximaler Auslenkung und der Kreisfrequenz $\omega$ lässt
sich die Auslenkung $y$ und die Geschwindigkeit $v$ in Abhängigkeit von der Zeit $t$ 
ausdrücken:

\begin{minipage}{.44\textwidth}
    \begin{equation}
        y = y_{max} \cdot \sin(\omega \cdot t)
    \end{equation}
\end{minipage}
\hspace{.1\textwidth}
\begin{minipage}{.44\textwidth}
    \begin{equation}
        v = \omega \cdot y_{max} \cdot \cos(\omega \cdot t)
    \end{equation}
\end{minipage}

Die Beschleunigung a in Abhängigkeit von der Zeit $t$ ist:

\begin{equation}
    a = -\omega^2 \cdot y_{max} \cdot \sin(\omega \cdot t) = - \omega^2 \cdot y
\end{equation}

Es zeigt sich, dass bei einer Sinusschwingung die Beschleunigung $a$ 
der Auslenkung $y$ proportional, ihr aber stets entgegengerichtet ist. 
Der Proportionalitätsfaktor ist negativ und hat einen Betrag gleich 
dem Quadrat der Kreisfrequenz \cite{zeitler_physik_2016}.

\subsubsection{Dynamik der Sinusschwingung}

Die kinematische Beschreibung der Sinusschwingung lässt die Frage offen, 
wodurch Frequenz und Amplitude der Schwingung bestimmt werden. Diese Frage
wird beantwortet, wenn wir die bei einer Schwingung auftretenden Kräfte und
Energien untersuchen und dabei das Grundgesetz der Dynamik und den
Energieerhaltungssatz anwenden \cite{zeitler_physik_2016}.

Wir betrachten dazu wiederum als Beispiel einen Masse-Feder-Schwinger, bei dem
ein Körper der Masse $m$ an einer Feder der Federkonstante $k$ hängt.

\begin{minipage}{.68\textwidth}
    Bei einer Auslenkung des Körpers um die Strecke $y$ aus seiner Gleichgewichtslage 
    wirkt die Federkraft $F = - ky$ als rücktreibende Kraft, die der 
    Auslenkung proportional und ihr entgegen gerichtet ist. (Beachten Sie, dass Sie hierbei
    die Gewichtskraft des Körpers nicht zu berücksichtigen brauchen, da sie bei y = 0
    durch eine gleich große Federkraft im Gleichgewicht gehalten wird. Die rück-
    treibende Kraft ist nur die Federkraft, die zusätzlich bei einer Auslenkung aus
    dieser Gleichgewichtslage auftritt.) Schwingungsfähige Systeme, für die ein 
    lineares Kraftgesetz gilt, sind harmonische Oszillatoren.
    Setzen wir die rücktreibende Kraft $- ky$ für $F$ in das Grundgesetz der Dynamik
    $F = ma$ ein, so erhalten wir daraus die Beschleunigung:

    \begin{equation}
        a = - \frac{k}{m} \cdot y = - \omega^2 \cdot y
    \end{equation}
\end{minipage}
\hspace{.05\textwidth}
\begin{minipage}{.35\textwidth}
    \begin{figure}[H]   
        \centering          %     left botm right top
        \includegraphics[clip, trim=0cm 0cm 0cm 0cm,width=\textwidth]{dynamik_sinsusschwingung.png} %.6 stands for 0.6 x textwidth
        \caption{Masse-Feder-Schwinger \cite{zeitler_physik_2016}}
        \label{fig:dynamiksinusschwingung}
    \end{figure}
\end{minipage}

Daraus kann die Eigenkreisfrequenz hergeleitet werden:

\begin{equbox}
    \begin{equation}
        \omega_0 = \sqrt{\frac{k}{m}}
    \end{equation}
\end{equbox}

\begin{bspbox}
    Wie gross sind Kreisfrequenz, Eigenfrequenz, Schwingungsdauer, 
    Amplitude, Maximalgeschwindigkeit und
    Maximalbeschleunigung der ungedämpften Schwingungen 
    eines Masse-Feder-Schwingers mit der Federkonstante 
    $25 Nm^{-1}$ und der Masse $100 g$? Die Schwingung
    wurde mit einer Energie von $20 mJ$ angeregt.

    Kreisfrequenz:

    \begin{equation*}
        \omega = \sqrt{\frac{k}{m}} = \sqrt{\frac{25Nm^{-1}}{0.1kg}} = 15.8s^{-1}
    \end{equation*}

    die Eigenfrequenz $f_0 = \frac{\omega}{2 \cdot \pi} = 2Hz$
    Es finden also 2.5 Schwingungen in einer Sekunde mit einer 
    Schwingungsdauer von $T = \frac{1}{f_0} = 0.4s$ statt. 
    Die Amplitude errechnet sich zu:

    \begin{equation*}
        y_{max} = \sqrt{\frac{2 \cdot E_s}{k}} = \sqrt{\frac{2 \cdot 0.02J}{24Nm^{-1}}} = 0.041m = 4.1cm
    \end{equation*}

    Der Körper schwingt mit der
    Maximalgeschwindigkeit 
    $v_{max} = \omega \cdot y_{max} = 15.8 s^{-1} \cdot 0.041 m
    = 0.65ms^{-1}$ durch die Nulllage und erfährt an den Umkehrpunkten 
    eine Maximalbeschleunigung $a_{max} = \omega^2 \cdot y_{max} 
    = (15.8s^{-1})^2 \cdot 0.041m = 10.2ms^{-2}$
\end{bspbox}

\subsubsection{Fadenpendel}

\begin{minipage}{.58\textwidth}
    Für das Fadenpendel gilt:

    \begin{equbox}
        \begin{equation}
            \omega_0 = \sqrt{\frac{g}{l}}
        \end{equation}
    \end{equbox}
    \vspace{10pt}

    Die Bewegung wird durch den Winkel beschrieben:

    \begin{equbox}
        \begin{align}
            \varphi(t) &= A \cdot \sin(\omega_0 \cdot t + \phi) \\
            \omega(t) &= A \cdot \omega_0 \cdot \cos(\omega_0 \cdot t + \phi) \\
            \alpha(t) &= - A \cdot \omega_0^2 \cdot \sin(\omega_0 \cdot t + \phi)
        \end{align}
    \end{equbox}

\end{minipage}
\hspace{.1\textwidth}
\begin{minipage}{.3\textwidth}
    \begin{figure}[H]    
        \centering          %     left botm right top
        \includegraphics[clip, trim=0cm 0cm 0cm 0cm,width=\textwidth]{fadenpendel_aufgabe.png}
        \label{fig:61skript}
    \end{figure}
\end{minipage}

\subsection{Freie gedämpfte Schwingungen}

\begin{minipage}{.68\textwidth}
    Eine freie ungedämpfte Schwingung schwingt unendlich lange, ihre Energie bleibt konstant.
    Reale Systeme in der Technik schwingen gedämpft, d.h. die Amplituden des freien
    gedämpften Schwingers nehmen ab, dem System wird Energie entzogen. Die Ursachen der
    Dämpfung können sehr unterschiedlich sein und reichen von innerer Reibung im Material
    über Lagerreibung (trocken oder geschmiert) bis zu viskoser Reibung und Luftwiderstand \cite{bartsch_dynamik_2023}.

\end{minipage}
\hspace{.05\textwidth}
\begin{minipage}{.35\textwidth}
    \begin{figure}[H]   
        \centering          %     left botm right top
        \includegraphics[clip, trim=0cm 0cm 0cm 0cm,width=\textwidth]{gedämpfte_schwingungen_gezeigt.png} %.6 stands for 0.6 x textwidth
        \caption{Freie gedämpfte Schwingung \cite{bartsch_dynamik_2023}}
        \label{fig:gedschwgeze}
    \end{figure}
\end{minipage}

Mit dem hier verwendeten Ansatz der Geschwindigkeitsproportionalen, viskosen Dämpfung
können viele Systeme in der Technik in guter Näherung beschrieben werden:

\begin{equation}
    F_D = d \cdot v = d \cdot \dot{x}
\end{equation}

Daraus kann die Abklingkonstante $\delta$ definiert werden:

\begin{equbox}
    \begin{equation}
        \delta = \frac{d}{2 \cdot m}
    \end{equation}
\end{equbox}

Oder mit dem Dämpfungsgrad $D$:

\begin{equbox}
    \begin{equation}
        D = \frac{\delta}{\omega_0}
    \end{equation}
\end{equbox}

Es sind die drei Fälle für eine gedämpfte Schwingung zu unterscheiden:

\begin{minipage}{.47\textwidth}
    \begin{infobox}
        \begin{itemize}
            \item Starke Dämpfung ($D > 1$)
            \item Kritische Dämpfung ($D = 1$)
            \item Schwache Dämpfung ($D < 1$)
        \end{itemize}
    \end{infobox}
    \vspace{10pt}

    Nachfolgend wird nur auf die schwache Dämpfung näher eingegangen, 
    da sie der relevanteste Fall für Schwingungen darstellt.
\end{minipage}
\hspace{.05\textwidth}
\begin{minipage}{.47\textwidth}
    \begin{figure}[H]   
        \centering          %     left botm right top
        \includegraphics[clip, trim=0cm 0cm 0cm 0cm,width=\textwidth]{dämpfungsarten.png} %.6 stands for 0.6 x textwidth
        \caption{Dämpfungsarten freie Schwingung \cite{noauthor_freie_nodate}}
        \label{fig:dämpfungsarten}
    \end{figure}
\end{minipage}

Das Verhältnis von zwei aufeinanderfolgenden Schwingungen ist ein Mass
für die Dämpfung. Mit Hilfe der Hüllkurve kann dieses Verhältnis $q$ 
berechnet werden (einhüllende, abnehmende Kurve abb. \ref{fig:gedschwgeze}).

\begin{equbox}
    \begin{align}
        \Lambda &= ln(q) \\
        D &= \frac{\Lambda}{\sqrt{4 \cdot \pi^2 + \Lambda^2}}
    \end{align}
\end{equbox}

Die schwache Dämpfung beeinflusst die Amplitude $A$ über die Zeit sowie die Eigenfrequenz $\omega_d$:

\begin{equbox}
    \begin{align}
        x(t) = A \cdot e^{- \delta \cdot t} \cdot \sin(\omega_d \cdot + \phi) \\
        \omega_d = \sqrt{\omega_0^2 - \delta^2}
    \end{align}
\end{equbox}

Mit Ausnahme spezieller Dämpfer (z.B. Fahrzeugachsen) betragen die Dämpfungsgrade $D$
von Maschinen und Strukturen meist nur wenige Prozente. Die Eigenfrequenz $\omega_d$ der freien
gedämpften Schwingung ist in diesem Fall nur unwesentlich kleiner als diejenige der
ungedämpften Schwingung $\omega_0$. Für $D = 20\%$ ändert sich die Eigenfrequenz um ca. 2\%.

Die Konstanten A und j sind wiederum aus den Anfangsbedingungen zu bestimmen.
Mit der Auslenkung $x_0$ bzw. der Geschwindigkeit $v_0$ zur Zeit $t = 0$ folgt:

\begin{equbox}
    \begin{align}
        A = \sqrt{x_0^2 + (\frac{v_0 + \delta \cdot x_0}{\omega_d})^2} \\
        \phi = \tan^{-1}(\frac{x_0 \cdot \omega_d}{v_0 + \delta \cdot x_0})
    \end{align}
\end{equbox}

\begin{bspbox}
    \begin{minipage}{.55\textwidth}
        Eine Masse $5 kg$ gleitet reibungsfrei auf einer schiefen Ebene. Die
        vorgespannten Federn haben eine Steifigkeit $c = 2 kNm^{-1}$, die
        Dämpfer $d = 25 Nsm^{-1}$. Die Masse wird zunächst 50 mm aus der
        Gleichgewichtslage nach oben geschoben und anschliessend mit
        einer Abwärtsgeschwindigkeit von $1.25 ms^{-1}$ losgelassen.
        Zu bestimmen sind:

        a) Die Periode der gedämpften Schwingung

        b) Die Lage der Masse in Funktion der Zeit

        Gravitation und Federvorspannung haben auf die Eigenfrequenz
        keinen Einfluss und können weggelassen werden.
        Federn und Dämpfer wirken parallel.

        \begin{align*}
            \omega_0 &= \sqrt{\frac{2 \cdot c}{m}} = \sqrt{\frac{2 \cdot 2000Nm^{-1}}{5kg}} = 28.28rads^{-1} \\
            \delta &= \frac{2 \cdot d}{2 \cdot m} = \frac{2 \cdot 25Nsm^{-1}}{2 \cdot 5kg} = 5rads^{-1} \\
            \omega_d &= \sqrt{\omega_0^2 - \delta^2} = \sqrt{(28.28rads^{-1})^2 - (5rads^{-1})^2} \\
            &= 27.84rads^{-1}
        \end{align*}

    \end{minipage}
    \hspace{.05\textwidth}
    \begin{minipage}{.38\textwidth}
        \begin{figure}[H]
            \centering
            \begin{subfigure}{\textwidth}
                \includegraphics[clip, trim=0cm 0cm 0cm 0cm, width=\textwidth]{beispiel_dämpfung.png}
                \label{fig:bspdämpfung}
            \end{subfigure}
        \end{figure}
    \end{minipage}

    Die Amplitude $A$ und der Phasenwinkel $\phi$ ergeben sich aus den Anfangsbedingungen wie folgt:

    \begin{align*}
        A &= \sqrt{x_0^2 + (\frac{v_0 + \delta \cdot x_0}{\omega_d})^2} 
        = \sqrt{(-0.05m)^2 + (\frac{1.25ms^{-1} 
        - 5rads^{-1} \cdot 0.05m}{27.84rads^{-1}})^2} = 0.0616m \\
        \phi &= \tan^{-1}(\frac{x_0 \cdot \omega_d}{v_0 + \delta \cdot x_0}) 
        = \tan^{-1}(\frac{-0.05m \cdot 27.84rads^{-1}}{1.25ms^{-1} - 5rads^{-1} \cdot 0.05m}) = -0.948rad \\
        x(t) &= 0.0616 \cdot e^{-5 \cdot t} \cdot \sin(27.84 \cdot t - 0.948)
    \end{align*}

\end{bspbox}

\section{Wellen}

Eine mechanische Welle ist die räumliche Ausbreitung einer Schwingung von Materie 
(mechanische Schwingung).Im Gegensatz zu anderen Wellentypen benötigen mechanische Wellen damit 
schwingfähige Medien (Gase, Fluide oder Festkörper), um sich ausbreiten zu können.

Mechanische Wellen transportieren hierbei nicht die Teilchen selbst, 
sondern wie jede andere Welle nur die Energie der anfänglichen Störung.
Zu den Beispielen zählen \cite{noauthor_mechanische_2023}:

\begin{minipage}{.38\textwidth}
    \begin{itemize}
        \item Wasserwellen
        \item Schallwellen
        \item seismische Wellen 
        \item allgemein Vibrationen
    \end{itemize}

\end{minipage}
\hspace{.05\textwidth}
\begin{minipage}{.55\textwidth}
    \begin{figure}[H]   
        \centering          %     left botm right top
        \includegraphics[clip, trim=0cm 0cm 0cm 0cm,width=\textwidth]{Espejo_(3207185886).jpg} %.6 stands for 0.6 x textwidth
        \caption{Wasserwellen als Beispiel für mechanische Oberflächenwellen \cite{noauthor_mechanische_2023}}
        \label{fig:wasserwellen}
    \end{figure}
\end{minipage}

Man unterscheidet grundsätzlich zwischen zwei verschiedenen Wellenarten:

\begin{enumerate}
    \item Querwelle (Wandern einer Querstörung)
    \item Längswelle (Wandern einer Längsstörung)
\end{enumerate}

\begin{figure}[H]
    \centering
    \begin{minipage}{0.4\textwidth}
        \includegraphics[width=\textwidth]{querwelle.png}
    \end{minipage}
    \begin{minipage}{0.4\textwidth}
        \includegraphics[width=\textwidth]{längswelle.png}
    \end{minipage}
    \caption{Beispiel für eine Querwelle (links) und eine Längswelle (rechts) \cite{boge_physik_2008}}
    \label{fig:quervslängs}
\end{figure}

\begin{minipage}{.48\textwidth}
    Wird eine Störung in den Wellenträger ein-
    gegeben, schwingen die Teilchen (Oszillatoren =
    Schwinger) am Ort, reißen aber durch ihre Kopp-
    lung die benachbarten Teilchen mit, sodass die
    Störung über den Wellenträger wandert. Jedes
    Teilchen gibt seine aufgenommene Energie fort-
    laufend an seine Nachbarteilchen ab, d.h. die
    Energie wird mit der fortschreitenden Störung
    transportiert.

\end{minipage}
\hspace{.05\textwidth}
\begin{minipage}{.45\textwidth}
    \begin{figure}[H]   
        \centering          %     left botm right top
        \includegraphics[clip, trim=0cm 0cm 0cm 0cm,width=\textwidth]{wellen_erklärung_oszillatoren.png} %.6 stands for 0.6 x textwidth
        \caption{Wandern der Störung \cite{boge_physik_2008}}
        \label{fig:wandernDerStörung}
    \end{figure}
\end{minipage}

Die Störung wandert, aber die Stoffteilchen
schwingen am Ort
Durch innere Reibung nimmt der Wellenträger
selbst Energie auf (es wird Reibarbeit verrichtet).

\subsection{Gleichung der harmonischen Welle}

Eine harmonische Welle entsteht, wenn die
einzelnen Schwinger (Oszillatoren) harmonische
Schwingungen ausführen.
Wellenlänge
$\lambda$ ist der Abstand zweier benachbarter Oszillatoren 1 und 2, die in gleicher Phase
schwingen, am zweckmässigsten gemessen als
Abstand zweier Wellenberge oder -täler.

Die Ausbreitungsgeschwindigkeit c der harmonischen Welle ist wie bei jeder 
gleichförmigen Bewegung der Quotient aus Weg- und
zugehörigem Zeitabschnitt.

\begin{minipage}{.48\textwidth}
    \begin{infobox}
        Die Welle legt während der Schwingungsdauer
        $T$ den Weg $\lambda$ zurück.
    \end{infobox}

    \begin{equbox}
        \begin{equation}
            c = \frac{\lambda}{T} = \lambda \cdot f
        \end{equation}
    \end{equbox}

\end{minipage}
\hspace{.05\textwidth}
\begin{minipage}{.45\textwidth}
    \begin{figure}[H]   
        \centering          %     left botm right top
        \includegraphics[clip, trim=0cm 0cm 0cm 0cm,width=\textwidth]{lambda.png} %.6 stands for 0.6 x textwidth
        \caption{$\lambda$ visualisiert \cite{boge_physik_2008}}
        \label{fig:lambda}
    \end{figure}
\end{minipage}

Um die harmonische Welle beschreiben zu
können, wird eine Gleichung entwickelt. Damit hat
man die Möglichkeiten:

\begin{itemize}
    \item für einen bestimmten gegebenen Zeitpunkt
    die Auslenkung $y$ eines jeden Wellenpunkts
    zu berechnen und
    \item für verschiedene Zeitpunkte die augenblickliche 
    Auslenkung $y$ eines einzelnen Wellenpunkts zu berechnen.
\end{itemize}

Da jeder Wellenpunkt der harmonischen Welle
eine harmonische Schwingung ausführt, wird von
der schon bekannten Gleichung der harmonischen Schwingung ausgegangen.

\begin{equation}
    y = A \cdot \sin(2 \cdot \pi \cdot \frac{t}{T})
\end{equation}

Diese Schwingungsgleichung enthält noch keine Grösse, die von der Ausbreitungsgeschwindigkeit
$c$ der Welle abhängig ist.

Die Schwingungsgleichung in der Form
$y = A \cdot \sin(2 \cdot \pi \cdot \frac{t}{T})$ enthält noch keine Grösse, die
von der Ausbreitungsgeschwindigkeit $c$ der Welle
abhängig ist.

\begin{minipage}{.58\textwidth}
    Beobachtet werden zwei Punkte 1 und 2 des
    Wellenträgers, die den Abstand $\Delta x$ voneinander
    haben. Wenn die Wellenbewegung Punkt 1 er-
    reicht, ist Punkt 2 noch in Ruhe (a). Die Welle
    läuft weiter und Punkt 1 schwingt zunächst nach
    oben (b) und dann wieder nach unten, bis die
    Welle auch den Punkt 2 erreicht hat (c). Punkt 2
    beginnt also seine Schwingung um einen 
    Zeitabschnitt $\Delta t$ später als Punkt 1. Auch während
    der weiteren Wellenbewegung schwingt Punkt 2,
    im gleichen Zeitabstand „hinter dem Punkt her“,
    d.h. in einem beliebigen Zeitpunkt $t$ hat Punkt 2
    dieselbe Auslenkung y, die Punkt 1 im Zeitpunkt
    $t - \Delta t$ hatte.

\end{minipage}
\hspace{.05\textwidth}
\begin{minipage}{.35\textwidth}
    \begin{figure}[H]   
        \centering          %     left botm right top
        \includegraphics[clip, trim=0cm 0cm 0cm 0cm,width=\textwidth]{herleitung_welle.png} %.6 stands for 0.6 x textwidth
        \caption{Herleitung \cite{boge_physik_2008}}
        \label{fig:herleitungWelle}
    \end{figure}
\end{minipage}

\begin{minipage}{.53\textwidth}

    Wenn Punkt 1 im Zeitpunkt t die Auslenkung
    $y = A \cdot \sin(2 \cdot \pi \cdot \frac{t}{T})$ hat, hat Punkt 2 erst die 
    Auslenkung $y = A \cdot \sin(2 \cdot \pi \cdot \frac{(t - \Delta t)}{T})$ erreicht.
    \vspace{8pt}
    
    Im Zeitabschnitt $\Delta t$ legt die Welle aber den
    Weg $\Delta x$ mit der Ausbreitungsgeschwindigkeit $c$
    zurück. Folglich ist $c = \frac{\Delta x}{\Delta t}$. Es ist aber auch
    $c = \frac{\lambda}{T}$.
    \vspace{8pt}
    
    Damit kann die Auslenkung $y$ des Punktes 2 (und
    damit jedes beliebigen Punktes) des Wellenträgers im Zeitpunkt $t$ bestimmt werden.

\end{minipage}
\hspace{.05\textwidth}
\begin{minipage}{.4\textwidth}
    \begin{align}
        c = \frac{\Delta x}{\Delta t} = \frac{\lambda}{T}; \frac{\Delta t}{T} = \frac{\Delta x}{\lambda} \\
        y = A \cdot \sin(2 \cdot \pi \cdot \frac{(t - \Delta t)}{T}) \\
        y = A \cdot \sin(2 \cdot \pi \cdot (\frac{t}{T} - \frac{\Delta t}{T}))
    \end{align}

    \begin{equbox}
        \begin{equation}
            y = A \cdot \sin(2 \cdot \pi \cdot (\frac{t}{T} - \frac{\Delta x}{\lambda}))
        \end{equation}
    \end{equbox}
\end{minipage}

